% part: intuitionistic-logic
% chapter: soundness-completeness
% section: lindenbaum

\documentclass[../../../include/open-logic-section]{subfiles}

\begin{document}

\olfileid{int}{sc}{lin}

\olsection{లిండెన్‌బామ్ ఉపప్రమేయం}

అంతఃప్రజ్ఞావాద తర్కపు సంపూర్ణతా సిద్ధాంతాన్ని
$\Gamma \Proves/ !A$ అని అనుకుని,
$\mSat{M}{\Gamma}$ మరియు~$\mSat/{M}{!A}$
అయ్యే నమూనాను నిర్మించడం ద్వారా నిరూపిస్తాం.

సాంప్రదాయిక తర్కంలో \tetoken{వ్యుత్పాద్యత}{!!{derivability}}
సంబంధాన్ని అవైరుధ్య భావనకు తగ్గించవచ్చు:
\tetoken{ఒక సూత్రం}{!!a{formula}}~$!A$ ఒక
\tetoken{సూత్రాల}{!!{formula}s} సమితి నుంచి
\tetoken{వ్యుత్పాదించదగినది}{!!{derivable}}
అయితే, అప్పుడే ఆ సమితికి~$!A$ నిరాకరణను
కలిపినది వైరుధ్యమైనది. అంతఃప్రజ్ఞావాద తర్కంలో
ఇది సాధ్యం కాదు. ఇందులో $\lnot!A$
వైరుధ్యమైనదైతే, $\Proves \lnot\lnot !A$
అని మాత్రమే వస్తుంది. $\lnot\lnot!A \lif !A$
సాధారణంగా అంతఃప్రజ్ఞావాద తర్కంలో నిలవదు
కాబట్టి~$\Proves
!A$ అని తేల్చలేం.

అందువల్ల నమూనా~$\mModel{M}$ను నిర్మించేటప్పుడు
\tetoken{సూత్రం}{!!{formula}}~$!A$ యొక్క
\tetoken{అవ్యుత్పాద్యతను}{!!{derivability}}
గమనిస్తూ ఉండాలి. అందుకే సాంప్రదాయిక
పద్ధతిలోలాగా నమూనా~$\mModel{M}$ నిర్మాణానికి
సంపూర్ణ సమితి $\Gamma^* \supseteq \Gamma$ను
ఉపయోగించలేం: ప్రతి సంపూర్ణ
సమితి $\Gamma^*$లో $\Gamma^* \Proves !A \lor \lnot !A$.

సంపూర్ణ సమితి~$\Gamma^*$కు బదులు, సూత్రాల
ప్రధాన సమితి అనే భావనను వాడుతాం:

\begin{defn}\ollabel{defn:prime}
  \tetoken{సూత్రాల}{!!{formula}s} సమితి~$\Gamma$
  \emph{ప్రధానం} అయితే, అప్పుడే
  \begin{enumerate}
  \item\ollabel{defn:prime1} $\Gamma$ అవైరుధ్యమైనది;
    అంటే $\Gamma \Proves/ \lfalse$;
  \item\ollabel{defn:prime2} $\Gamma \Proves !A$
    అయితే $!A \in \Gamma$; మరియు
  \item\ollabel{defn:prime3} $!A \lor !B \in \Gamma$
    అయితే $!A \in
    \Gamma$ లేదా $!B \in \Gamma$.
  \end{enumerate}
\end{defn}

\begin{lem}[లిండెన్‌బామ్ ఉపప్రమేయం]
  \ollabel{lem:lindenbaum} $\Gamma \Proves/ !A$
  అయితే, $\Gamma^* \supseteq \Gamma$,
  $\Gamma^*$ ప్రధానం, $\Gamma^* \Proves/ !A$
అయ్యే విస్తరణ ఉంది.
\end{lem}

\begin{proof}
  $!B \lor !C$ రూపంలోని అన్ని
  \tetoken{సూత్రాల}{!!{formula}s}ను
  $!B_1 \lor !C_1$, $!B_2 \lor !C_2$, \dots,
  అని వరుసపెడదాం. \tetoken{సూత్రాల}{!!{formula}s}
  సమితుల పెరుగుతున్న క్రమం~$\Gamma_n$ను
  నిర్వచిస్తాం. ప్రతి $\Gamma_{n+1}$,
  $\Gamma_n$కు ఒక కొత్త \tetoken{సూత్రం}{!!{formula}}
  చేర్చినదిగా ఉంటుంది. $\Gamma^*$ అనేది
  అన్ని~$\Gamma_n$ల సమ్మేళనం. కొత్త
  \tetoken{సూత్రాలను}{!!{formula}s} $\Gamma^*$
  ప్రధానంగా ఉండి, ఇంకా $\Gamma^* \Proves/ !A$
  నిలిచేలా ఎంచుకుంటాం. అంటే ప్రతి దశలో కింది
  షరతులు గల తొలి వియోజనం~$!B_i \lor !C_i$ను
  కనుగొనాలి:
  \begin{enumerate}
  \item\ollabel{gamma-1} $\Gamma_n \Proves !B_i \lor !C_i$
  \item\ollabel{gamma-2} $!B_i \notin \Gamma_n$ మరియు
    $!C_i \notin \Gamma_n$
  \end{enumerate}
  $\Gamma_n \cup \{!B_i\} \Proves/ !A$ అయితే
  $!B_i$ను, లేకపోతే $!C_i$ను $\Gamma_n$కు
  చేర్చుతాం. ఇది పనిచేస్తుందని చూపాలి.
  ప్రస్తుతానికి \olref{gamma-1},
  \olref{gamma-2} రెండూ నెరవేర్చే అతి చిన్న
  $i$ను $i(n)$గా నిర్వచిద్దాం.

  $\Gamma_0 = \Gamma$ అని, అలాగే
  \[
  \Gamma_{n+1} = \begin{cases}
    \Gamma_n \cup \{!B_{i(n)}\} &
    \text{$\Gamma_n \cup \{!B_{i(n)}\} \Proves/ !A$ అయితే} \\
    \Gamma_n \cup \{!C_{i(n)}\} & \text{లేకపోతే}
    \end{cases}
  \]
  అని నిర్వచిద్దాం. $i(n)$ నిర్వచితం కాకపోతే,
  అంటే $\Gamma_n \Proves !B \lor !C$ అయినప్పుడల్లా
  $!B \in \Gamma_n$ లేదా $!C \in \Gamma_n$
  అయితే, $\Gamma_{n+1}
  = \Gamma_n$గా తీసుకుంటాం. ఇప్పుడు
  $\Gamma^* = \bigcup_{n=0}^\infty \Gamma_n$.

  ముందుగా అన్ని $n$లకు $\Gamma_n \Proves/ !A$
  అని చూపుతాం. $n$పై ఆగమన పద్ధతిలో సాగుదాం.
  $n = 0$కు ఇది సిద్ధాంత పరికల్పన
  $\Gamma \Proves/ !A$ వల్ల నిలుస్తుంది.
  $n>0$ అయితే $\Gamma_n \Proves/ !A$ నుంచి
  $\Gamma_{n+1} \Proves/ !A$ అని చూపాలి.
  $i(n)$ నిర్వచితం కాకపోతే $\Gamma_{n+1} =
  \Gamma_n$; నిరూపించాల్సిందేమీ లేదు.
  కాబట్టి $i(n)$ నిర్వచితమని అనుకుందాం.
  సరళత కోసం $i =
  i(n)$ అని వ్రాద్దాం.

  వాదన ప్రతివిపర్యయాన్ని నిరూపిస్తాం.
  $\Gamma_{n+1}
  \Proves !A$ అనుకుందాం. నిర్మాణం ప్రకారం
  $\Gamma_n \cup \{!B_i\} \Proves/ !A$
  అయితే $\Gamma_{n+1} = \Gamma_n \cup
  \{!B_i\}$; లేకపోతే $\Gamma_{n+1} =
  \Gamma_n \cup \{!C_i\}$. మొదటి కేసు
  అసాధ్యం, ఎందుకంటే అప్పుడు
  $\Gamma_{n+1} \Proves/ !A$. అందువల్ల
  $\Gamma_n \cup
  \{!B_i\} \Proves !A$ మరియు
  $\Gamma_{n+1} = \Gamma_n \cup \{!C_i\}$.
  $i(n)$ నిర్వచనం వల్ల $\Gamma _n \Proves
  !B_i \lor !C_i$. మనకు $\Gamma_n \cup
  \{!B_i\} \Proves !A$ కూడా ఉంది;
  $\Gamma_{n+1} = \Gamma_n \cup \{!C_i\}
  \Proves !A$ కూడా ఉంది. కాబట్టి
  కోరినట్లుగా $\Gamma_n
  \Proves !A$.

  $\Gamma^* \Proves !A$ అయితే, $\Gamma'
  \subseteq \Gamma^*$ మరియు $\Gamma'
  \Proves !A$ అయ్యే ఏదో పరిమిత
  ఉపసమితి ఉంటుంది. ప్రతి $!D \in
  \Gamma'$ ఏదో~$i$కు $\Gamma_i$లో ఉండాలి.
  ఆ ఉపసమితి ఖాళీ కానప్పుడు ఈ సూచికల్లో
  అతి పెద్దదాన్ని $n$గా తీసుకుందాం;
  ఖాళీ అయితే మొదటి దశ సూచికనే తీసుకుందాం.
  \sourcecorrection{OLTEINTLIN-001}{మూలం పరిమిత $\Gamma'$ ఖాళీగా ఉండే అవకాశాన్ని పరిగణించకుండా దాని సభ్యుల సూచికల్లో గరిష్ఠాన్ని తీసుకుంది; ఖాళీ సందర్భంలో మొదటి దశను తీసుకునేలా స్పష్టం చేశాం.}
  $i \le n$ అయితే $\Gamma_i \subseteq
  \Gamma_n$ కాబట్టి $\Gamma'
  \subseteq \Gamma_n$. అయితే
  $\Gamma_n \Proves !A$ వస్తుంది;
  ఇది పైన నిరూపించిన $\Gamma_n \Proves/ !A$కు
  విరుద్ధం.

  చివరగా $\Gamma^*$ ప్రధానమని, అంటే
  \olref{defn:prime1}, \olref{defn:prime2},
  \olref{defn:prime3} షరతులను
  \olref{defn:prime}లో చూపినట్లుగా
  నెరవేర్చుతుందని చూపుతాం.

  మొదట $\Gamma^* \Proves/ !A$ కాబట్టి
  $\Gamma^*$ అవైరుధ్యమైనది; అందువల్ల
  \olref{defn:prime1} నిలుస్తుంది.

  ఇప్పుడు $\Gamma^* \Proves !B \lor !C$
  అయితే $!B
  \in \Gamma^*$ లేదా $!C \in \Gamma^*$
  అని చూపుతాం. దీనివల్ల \olref{defn:prime3}
  నిలుస్తుంది; ఎందుకంటే $!B \lor !C \in
  \Gamma^*$ అయితే $\Gamma^* \Proves !B
  \lor !C$ కూడా. $\Gamma^* \Proves
  !B \lor !C$, కానీ $!B \notin \Gamma^*$,
  $!C \notin \Gamma^*$ అని అనుకుందాం.
  $\Gamma^* \Proves !B \lor !C$ కాబట్టి
  ఏదో~$n$కు $\Gamma_n \Proves !B \lor !C$.
  అన్ని వియోజనాల వరుసలో $!B \lor
  !C$ అనేది, ఉదాహరణకు, $!B_j \lor !C_j$గా
  ఉంటుంది. $!B_j \lor !C_j$ అనేది $i(n)$ నిర్వచనంలోని
  షరతులు నెరవేర్చుతుంది: $\Gamma_n
  \Proves !B_j
  \lor !C_j$, కాగా $!B_j \notin \Gamma_n$,
  $!C_j \notin \Gamma_n$. ఆ సూచిక కంటే ముందు
  ఉన్న సూచికలు పరిమిత సంఖ్యలోనే ఉన్నాయి;
  వాటిలో ఏ $!B_i \lor !C_i$ తొలి అర్హ
  వియోజనంగా ఎంపికైనా, దాని $!B_i$ లేదా
  $!C_i$ వియోజ్యాన్ని చేర్చిన
  తరువాత అది మళ్లీ అర్హం కాదు. కాబట్టి
  కొంత దశ~$m$లో $j = i(m)$ అవుతుంది.
  \sourcecorrection{OLTEINTLIN-002}{మూలం ప్రతి దశలో మొత్తం అర్హ వియోజనాల సంఖ్య తగ్గుతుందని పేర్కొంది; కొత్త వ్యుత్పాద్యతల వల్ల అది నిర్ధారణ కాదు. స్థిర $j$కి ముందు పరిమిత సూచికల్లో ప్రతిదీ ఒక్కసారే ఎంపికవుతుందనే సరైన కారణాన్ని వాడాం.}
  కానీ అప్పుడు $!B \in \Gamma_{m+1}$
  లేదా $!C \in \Gamma_{m+1}$; ఇది
  $!B \notin \Gamma^*$,
  $!C \notin \Gamma^*$ అనే పరికల్పనకు
  విరుద్ధం.

  ఇప్పుడు $\Gamma^* \Proves !B$ అనుకుందాం.
  అప్పుడు $\Gamma^* \Proves !B \lor
  !B$. $\Gamma^* \Proves !B \lor !B$
  అయితే $!B \in \Gamma^*$ అని ఇప్పుడే
  నిరూపించాం. కాబట్టి $\Gamma^*$
  \olref{defn:prime2}ను
  \olref{defn:prime}లో చెప్పినట్లుగా నెరవేర్చుతుంది.
\end{proof}

\begin{prob}
$\Gamma \Proves/ \bot$ అయితే $\Gamma$
సాంప్రదాయిక తర్కంలో అవైరుధ్యమైనదని, అంటే
ఆ~$\Gamma$లోని అన్ని సూత్రాలను సత్యం చేసే
విలువకేటాయింపు ఉందని చూపండి.
\end{prob}

\end{document}
